\subsection{Shift-equivalence universality}

\begin{theorem}[Universality class $=$ fixed-dimension shift-equivalence class]\label{thm:universality-se-class}\lean{Obsmetrologia.AperiodicOrder.Tiling.Universality.universality_class_se_invariance}\leanok\uses{thm:mld-refines-cohequiv, thm:clinv-dichotomy}
At a fixed dimension $n$, two primitive non-singular step-$n$ substitution
matrices have the same universality class iff they are shift equivalent (the
class is $\mathrm{Quot}\,\mathrm{SE}$), and the cohomological classifying
invariant $\mathrm{cl\text{-}inv}$ is SE-invariant (rules sharing the
classifying coordinate carry the same class)
(\texttt{universality\_class\_se\_invariance}).
\end{theorem}

\begin{theorem}[Cross-dimensional shift equivalence]\label{thm:universality-cross-dim-se}\notready\uses{thm:universality-se-class}
\emph{Open frontier, stated but not formalized.} The general universality
relation is shift equivalence across dimensions, where the strong shift
equivalence may change the matrix rank (Williams' classification of subshifts
of finite type up to conjugacy). The core models SE only at a fixed $n$
(\cref{thm:universality-se-class}); the rank-changing relation and the
decidability status of the Williams problem are the open frontier. Red until
cross-dimension SE is modeled.
\end{theorem}
